\chapter*{Lời mở đầu}
\addcontentsline{toc}{chapter}{Lời mở đầu}

\begin{tcolorbox}[
	enhanced,
	colback=white,
	colframe=goldaccent!70!chapterline,
	boxrule=0.7pt,
	arc=2mm,
	left=3mm,
	right=3mm,
	top=2.2mm,
	bottom=2mm,
	borderline west={1.5mm}{0pt}{tealaccent!80!black}
]
{\small\bfseries\color{chapterline} Tuyên ngôn sử dụng nhanh\par}
\vspace{0.3em}
{\large\bfseries\color{titlecolor} Mở đúng chương, đọc đúng lúc, cười xong phải tỉnh.\par}
\vspace{0.28em}
{\footnotesize\itshape\color{noteink} Open the right chapter, read at the right moment, then leave the laugh with a useful sting.\par}
\end{tcolorbox}

Cuốn này được chỉnh lại theo đúng ý đồ ban đầu: \textbf{một sách thơ song ngữ đủ tám mươi chương}, ưu tiên nhịp đọc miệng gọn, độ cà khịa sắc, và cảm giác dùng thật được trên lớp. Mỗi chương chạm vào một thói quen học đường rất cụ thể: né bài, xin điểm, giấu điện thoại, mơ mộng, trì hoãn, nói nhiều, hứa đổi đời rồi ngủ tiếp. Vì vậy người dạy có thể mở trúng đúng tật, đọc đúng một bài, và chốt luôn ý giáo dục mà không phải diễn giải dài dòng.

Tinh thần của bản này là \textbf{đâm đúng tật chứ không đâm cho sướng miệng}. Câu thơ phải có lực, phải khiến học sinh bật cười vì thấy mình ở trong đó; nhưng sau tiếng cười vẫn phải còn một mũi kim: học đừng đối phó, hứa đừng lấy lệ, hiểu rồi hãy chép, trượt rồi đừng diễn quá lâu, quyết tâm thì phải bền hơn một đêm.

\begin{tcolorbox}[
	enhanced,
	colback=softgray,
	colframe=softline,
	boxrule=0.5pt,
	arc=2mm,
	left=2.6mm,
	right=2.6mm,
	top=1.6mm,
	bottom=1.5mm
]
{\small\bfseries\color{titlecolor} Dùng cuốn này theo ba cách\par}
\vspace{0.25em}
\begin{itemize}
\item Đọc mở tiết: lấy một bài ngắn để kéo sự chú ý, rồi vào bài chính.
\item Đọc chốt tiết: dùng một bài thơ để kết luận đúng cái tật lớp vừa phạm.
\item Đọc chữa thói quen: giao cả chương cho lớp tự soi, tự cười, tự sửa.
\end{itemize}
\end{tcolorbox}

Toàn bộ nội dung vẫn lấy \textbf{thất ngôn tứ tuyệt} làm trục vì đây là thể phù hợp nhất cho kiểu đọc ngẫu hứng. Bốn câu Việt đủ ngắn để lên nhịp, dừng nhịp, bẻ giọng; phần tiếng Anh đi kèm giữ tinh thần bài gốc, dùng như một cú nhấn song ngữ chứ không kéo tiết học đi chỗ khác. Từ bản này trở đi, mục lục cũng hiện cả tên từng bài thơ để người đọc lật đúng mục tiêu nhanh hơn.

Đích cuối cùng vẫn rất rõ: \textbf{đọc cho vui chỉ là bước một}. Đọc mà cả lớp cười xong thấy nhột, thấy đúng, thấy nên tự mở vở ra thì bài thơ mới làm tròn việc của nó.